\subsection{Product measurements and the full-AA reduction}
\label{quantum-clusters:measurement}

\begin{corollary}[Full AA for measured local quantum Gibbs states]
\label{cor:quantum-clusters-full-aa}
Let $H=\sum_A h_A$ have $|A|\le k$ and
$\Lambda=\max_i\sum_{A\ni i}\|h_A\|$, including one-site terms.
If
\[
 \beta\Lambda\le\frac{e^{-k(2+\frac12\log2)}}{32ek},
\]
then the outcomes of any binary product POVM on
$e^{-\beta H}/\operatorname{Tr}(e^{-\beta H})$ satisfy full AA.
Every $T$-query acceptance function has a deterministic classical
simulator with mean error $\epsilon$ using
$O(T^{18}\epsilon^{-28})$ individual-outcome queries, with an
absolute implied constant. The same bound holds after arbitrary
positive-probability outcome pinnings. There is no degree,
commutativity, or local Hilbert-space dimension restriction.
\end{corollary}

\begin{proof}
At each site let \(E_i(+1),E_i(-1)\) be a binary POVM, so
\(E_i(+1)+E_i(-1)=I\).  The distribution obtained by applying these
local measurements to the Gibbs state is
\[
 \mu_\beta(x)=
 \frac{\operatorname{Tr}[(\bigotimes_i E_i(x_i))e^{-\beta H}]}
      {\operatorname{Tr}(e^{-\beta H})}.
\]
If all outcome effects have positive trace, set
\[
 \tau_i(x_i)=\operatorname{Tr}E_i(x_i),\qquad
 \varphi_{i,x_i}(B)=
     \frac{\operatorname{Tr}[E_i(x_i)B]}{\tau_i(x_i)}.
\]
These are positive normalized functionals, and the outcome law becomes
\[
 \mu_\beta(x)\ \propto\
 \prod_i\tau_i(x_i)\,
 \exp\!\left(\sum_{S\subseteq[n]}
                  \widehat V_\beta(S)\chi_S(x)\right).
 \tag{QMEAS}
\]
The product prefactor adds only one-spin fields and a constant to the
logarithmic potential.  Consequently it leaves the weighted quantity
in (QW) unchanged, even if the measurement biases are very large.
An effect with zero trace is zero, so its outcome is impossible; omit
that outcome coordinate or regard it as deterministically pinned.
The same proof applies to the remaining cube, with the fixed local
functional at the removed coordinate.  Projective measurements in any
product basis are a special case.

Thus Theorem~\ref{thm:quantum-clusters-potential} reduces these measured
noncommuting Gibbs laws to classical laws with arbitrary fields and
\[
 W=\max_i\sum_{S\ni i}(|S|-1)|J_S|\le\frac1{31k}<\frac18,
 \qquad J_S=\widehat V_\beta(S)\quad (|S|\ge2).
\]
Apply Theorem~\ref{thm:resolution-dynamics-arbitrary-degree} to this
classical potential, represented at inverse temperature one. This gives
the stated query bound. No appeal to correlation decay or to classical
simulability of the underlying quantum Gibbs state is needed.
It also preserves hereditary versions: fixing measured coordinates
turns a monomial \(\chi_S\) into a signed monomial on a subset of
\(S\), and the triangle inequality shows that the remaining weighted
incident norm cannot increase.
\end{proof}

For example, the transverse-field Ising Hamiltonian
$H=-\sum_{ij}J_{ij}Z_iZ_j-\sum_i h_iZ_i-\sum_i g_iX_i$ is covered
with $k=2$ and
$\Lambda=\max_i(\sum_j|J_{ij}|+|h_i|+|g_i|)$.
This includes arbitrary graphs and signs, for a fixed product
measurement. It establishes a sufficient high-temperature interval,
not an exact transition temperature.

\subsection{Scope and verification record}
\label{quantum-clusters:provenance}

The connected-multiset representation, the degree-independent rooted
tree estimate (QR), and the Fourier-weight conversion (QW) are derived
in this section.  The imported input is the abstract-polymer convergence
and pinned bound of Koteck\'y and Preiss, in the explicitly checked
formulation of Fern\'andez and Procacci,
\emph{Cluster expansion for abstract polymer models: New bounds from an
old approach},
\href{https://arxiv.org/abs/math-ph/0605041}{arXiv:math-ph/0605041},
Section 2, equations (2.7), (2.12)--(2.15).  That source provides the
countable-polymer and self-incompatibility conventions used here.
The original result is Koteck\'y--Preiss,
\emph{Communications in Mathematical Physics} 103 (1986), 491--498,
\href{https://doi.org/10.1007/BF01211762}{doi:10.1007/BF01211762}.

This argument concerns a fixed product measurement whose entire outcome
string is the classical query input.  The bound includes physical
one-site Hamiltonian terms in \(\Lambda\); it does not assert that
unbounded noncommuting physical fields can be ignored.  It places no
computational-time bound on finding the effective coefficients or on
preparing the Gibbs state.  The temperature constant is deliberately
conservative.  No claim of novelty relative to the cluster-expansion
literature is made.
